\begin{lemma}[The distributional gradient of an affinely normalized function]
Let $p\in\mathbb{R}$, let $\Omega\subseteq\mathrm{Plane}$, let $u\colon\mathrm{Plane}\to
\mathbb{R}$ and $g\colon\mathrm{Plane}\to\mathrm{Plane}$, let $x_0\in\mathrm{Plane}$, let
$r,M,c\in\mathbb{R}$ with $r>0$, and suppose $u$ is $p$-harmonic on $\Omega$ with weak gradient
$g$ in the sense of \noderef{IsPHarm}.  Write $T(x):=x_0+r\,x$ and
$\Omega':=T^{-1}\Omega=\{x:x_0+r\,x\in\Omega\}$.  Then for every
$\varphi\colon\mathrm{Plane}\to\mathbb{R}$ of class $C^{\infty}$ with compact support and
$\operatorname{supp}\varphi\subseteq\Omega'$,
\[
  \int_{\Omega'}\frac{u(T(x))-c}{M}\;D\varphi(x)\,dx
  =-\int_{\Omega'}\varphi(x)\;\frac{r}{M}\,g(T(x))\,dx .
\]

That is: the normalized function $x\mapsto\bigl(u(T(x))-c\bigr)/M$ has
$x\mapsto(r/M)\,g(T(x))$ as a weak gradient on $\Omega'$.  This is the distributional half of
affine normalization and one of the bridge facts the paper's blow-up layer needs.  The iteration
rescales its solution at every stage, and a rescaled function is a weak solution of anything only
once its weak gradient has been identified; that identification is this statement.

Two features are worth stating.  The additive constant $c$ is carried because the iteration
subtracts one --- the case $w(x,y)=\bigl(u(x+y/4)-u(x)\bigr)/M_0$ --- and it DISAPPEARS from the
right-hand side.  That is not bookkeeping but the one genuinely analytic step of the proof: the
mean of the gradient of a compactly supported test function vanishes, so the $c$ contribution
integrates to zero.  And no hypothesis $M>0$ is needed: $M$ enters only as a formal divisor,
linearly on both sides, so the identity holds for every $M$, including $M=0$, where it reads
$0=0$.
\end{lemma}
\begin{proof}
Note $r\ne0$ since $r>0$, so $T$ is invertible with $T^{-1}(y)=S(y):=r^{-1}(y-x_0)$.  Let
$\varphi$ be as given and put $\psi:=\varphi\circ S$.

\medskip\noindent\emph{Step 1: the pushed-forward test function.}
By \noderef{AffineTestFunctionPullback}, applied with this $x_0$, this $r>0$, the set $\Omega$ and
this $\varphi$ --- whose hypothesis on $\operatorname{supp}\varphi$ is exactly
$\operatorname{supp}\varphi\subseteq T^{-1}\Omega$, as assumed --- the function $\psi$ is of class
$C^{\infty}$, has compact support, satisfies $\operatorname{supp}\psi\subseteq\Omega$, and obeys
\begin{equation}\label{eq:nwg-chain}
  D\varphi(x)=r\,D\psi\bigl(T(x)\bigr)\qquad\text{for every }x .
\end{equation}
So $\psi$ is an admissible test function for $\Omega$ itself.  Also $\psi(T(x))=\varphi(x)$ for
every $x$, since $S(T(x))=x$.

\medskip\noindent\emph{Step 2: localizing both integrals to the whole plane.}
By \noderef{VanishingOffTsupport}, $\varphi$ and $D\varphi$ vanish off
$\operatorname{supp}\varphi$, and $\psi$ and $D\psi$ vanish off $\operatorname{supp}\psi$.  Since
$\operatorname{supp}\varphi\subseteq\Omega'$, both integrands on the left and right of the
assertion vanish off $\Omega'$, so each set integral over $\Omega'$ equals the corresponding
integral over all of $\mathrm{Plane}$.  Likewise $\operatorname{supp}\psi\subseteq\Omega$ lets us
replace $\int_\Omega$ by $\int_{\mathrm{Plane}}$ in the two integrals appearing below.  All four
replacements are instances of the clause of the integrability proposition that a function
vanishing off a subset has the same integral over the subset as over the ambient space.

\medskip\noindent\emph{Step 3: the weak-gradient identity for $u$, applied to $\psi$.}
Since $\psi$ is admissible for $\Omega$, the weak-gradient clause of
$p$-harmonicity for $u$ gives
\begin{equation}\label{eq:nwg-key}
  \int_{\mathrm{Plane}}u(y)\,D\psi(y)\,dy=-\int_{\mathrm{Plane}}\psi(y)\,g(y)\,dy ,
\end{equation}
after the replacements of Step 2.  Both integrals converge: $D\psi$ is continuous with compact
support, hence integrable, and $u\,D\psi$ is continuous on the compact $\operatorname{supp}\psi$
--- $u$ being continuous on $\Omega\supseteq\operatorname{supp}\psi$ --- and vanishes off it.

\medskip\noindent\emph{Step 4: the constant drops out.}
By \noderef{TestFunctionGradientIntegral} applied to $\psi$, which is of class $C^{1}$ with
compact support,
\[
  \int_{\mathrm{Plane}}D\psi(y)\,dy=0 .
\]
Hence, $c\,D\psi$ being integrable as a scalar multiple of an integrable function,
\begin{equation}\label{eq:nwg-split}
  \int_{\mathrm{Plane}}\bigl(u(y)\,D\psi(y)-c\,D\psi(y)\bigr)dy
  =\int_{\mathrm{Plane}}u(y)\,D\psi(y)\,dy-c\int_{\mathrm{Plane}}D\psi(y)\,dy
  =\int_{\mathrm{Plane}}u(y)\,D\psi(y)\,dy .
\end{equation}
This is where $c$ disappears, and it is the only place any hypothesis about $\psi$ beyond
admissibility is used.

\medskip\noindent\emph{Step 5: change of variables, and the cancellation of the Jacobian.}
By the affine change of variables for Lebesgue measure on the plane, for every integrable
$F\colon\mathrm{Plane}\to\mathrm{Plane}$,
\begin{equation}\label{eq:nwg-cov}
  \int_{\mathrm{Plane}}F\bigl(T(x)\bigr)dx
  =\bigl\lvert r^{-n}\bigr\rvert\int_{\mathrm{Plane}}F(y)\,dy ,
  \qquad n:=\dim\mathrm{Plane},
\end{equation}
the factor being the Jacobian of the dilation.  Now compute, using \eqref{eq:nwg-chain} and
$\psi(T(x))=\varphi(x)$ to rewrite each integrand, then \eqref{eq:nwg-cov} twice, then
\eqref{eq:nwg-split} and \eqref{eq:nwg-key}:
\[
\begin{aligned}
  \int_{\mathrm{Plane}}\frac{u(T(x))-c}{M}D\varphi(x)\,dx
  &=\frac{r}{M}\int_{\mathrm{Plane}}
      \bigl(u(T(x))\,D\psi(T(x))-c\,D\psi(T(x))\bigr)dx\\
  &=\frac{r}{M}\,\bigl\lvert r^{-n}\bigr\rvert
      \int_{\mathrm{Plane}}\bigl(u(y)\,D\psi(y)-c\,D\psi(y)\bigr)dy\\
  &=\frac{r}{M}\,\bigl\lvert r^{-n}\bigr\rvert
      \int_{\mathrm{Plane}}u(y)\,D\psi(y)\,dy
   =-\frac{r}{M}\,\bigl\lvert r^{-n}\bigr\rvert
      \int_{\mathrm{Plane}}\psi(y)\,g(y)\,dy\\
  &=-\frac{r}{M}\int_{\mathrm{Plane}}\psi\bigl(T(x)\bigr)g\bigl(T(x)\bigr)dx
   =-\int_{\mathrm{Plane}}\varphi(x)\,\frac{r}{M}\,g\bigl(T(x)\bigr)dx .
\end{aligned}
\]
The first line pulls the constant $r/M$ out by \eqref{eq:nwg-chain} and linearity; the second and
the fifth are \eqref{eq:nwg-cov}, applied to $y\mapsto u(y)D\psi(y)-cD\psi(y)$ and to
$y\mapsto\psi(y)g(y)$ respectively.  THE JACOBIAN FACTOR APPEARS ON BOTH SIDES AND CANCELS
between those two applications, which is why no dimensional constant survives into the statement.
By Step 2 the two outer integrals equal the corresponding integrals over $\Omega'$, which is the
assertion.  As $\varphi$ was an arbitrary admissible test function for $\Omega'$, the lemma is
proved.
\end{proof}
